\chapter{Ten books and two hands}

\section{What the books are}

On the inside front cover of the first of the six ledgers, in a small upright running hand, Josephine Nivison Hopper wrote what the books were for:

\begin{leafquote}
Recorded by Jo N Hopper (Mrs. Edward Hopper) at time each work is finished, or before it leaves studio. Drawings in the 3 books done by Edward Hopper.
\end{leafquote}

\noindent The second volume repeats the statement in a shorter form, in the same hand, inside a hand-drawn frame with bracketed corners that also carries the address, \q{Edward Hopper / 3 Washington Square North / New York City / --- / His Work. / Vol. II.}: \q{Recorded by Jo N. Hopper (Mrs. Edward Hopper.) before each work leaves studio. Drawings by Edward Hopper.} (\lfu{Book II}{inside cover}{17362}). The word \q{each} is squeezed in above the line, between \q{before} and \q{work}, which is the first of several thousand corrections she would make to her own record over forty years.

The statement is worth pausing on because it is a claim about method, and the whole documentary value of the archive turns on whether it is true. It says that the record was made at the moment the work left the studio, not reconstructed afterwards; that the drawings are his and everything else is hers; and that there are three books of that kind. All three claims are borne out by the leaves, and all three need qualifying.

There are six volumes at the Whitney, accessioned as 96.208 to 96.213, and they are not all of one kind. Three of them are the \q{3 books} of the inside cover: Book I (96.208), which begins with the etchings one plate to a leaf and then, from leaf 45, turns into running lists; Book II (96.209), the paintings one to an opening, with Edward Hopper's ink record sketch and Josephine Hopper's description opposite; and Book III (96.210), the last of the work books, which runs to his death in 1967 and carries the latest hand in the archive. The fourth, Book IV (96.211), is something else altogether: a stationer's pocket memorandum book, seven and a half inches by four and three-quarters, with no sketches and no anecdote, kept without a gap from 15 November 1913 to 23 March 1967. The fifth, Book V (96.212), is sixteen sheets of drawings and etchings, 1953 to 1963, which declares on its own index leaf, \q{This book for drawings, not water colors.} (\lfu{Book V}{index}{16748}). And the sixth, Dealers/Etchings (96.213), is a cash book whose front half was used for something else, so that its leaves begin at 51, and which is indexed, as the transcription's header puts it, \q{by whose hands the work was in} rather than by which work it was.

The three work books were the gift of Lloyd Goodrich; Book IV was a purchase, and the dealers' book a purchase with funds from an anonymous donor. That the friend who organised the retrospectives of 1950 and 1964 should have owned the record books is not surprising, and it is part of the history of the archive that the books were consulted, and their figures copied, by the museum long before they came to it. Josephine Hopper's own diary records, on 24 June 1948, that \q{Miss Irven bragged that she had all the data from my ledger, so they can go right ahead \& put their figures on everything} (PAAM typescript, Part V); Rosalind Irvine of the Whitney was preparing the 1950 retrospective, and the ledger she had the data from is Book II.

Beside the six ledgers, the Whitney's Sanborn Hopper Archive holds, by the museum's own count, ninety notebooks of Josephine Hopper's. Four of them have been digitised, and those four are read here: the notebook she titled \q{Garrulities + Grouch} and subtitled \q{When artists meet} (1924--1950, 42 sheets); two notebooks of 1947 on the fight with New York University over the house at 3 Washington Square North, one headed \q{Re: 3 Wash Sq} (51 sheets) and the other \q{Battle of Wash Sq.} (8 sheets); and a black two-ring notebook of about 1952 to 1961 (68 sheets) that is not a diary at all but a book of lists: frames and their makers, insurance, paper by the pound, addresses, mailing lists, a recipe, and two chronicles of the years, one kept forwards and one backwards. The remaining diaries, twenty-two of them at the Provincetown Art Association and Museum and the rest at the Whitney, are not in this corpus and are used here only through Madeleine Larson's typescript of the Provincetown volumes, which the Timeline of the edition cites page by page.

\section{Whose hand}

The division of labour the inside cover announces is real and it is visible on every leaf. Edward Hopper's contribution to the record books is the drawing and its caption: a small pen sketch of the finished picture, ruled into a box, in bold strokes and close cross-hatching, and above or beside it the title and the size in drawn lettering. The size is the one figure in the archive that he wrote himself, and it is also the only quantity in it that measures an object rather than reporting an event. Everything else is hers: the description, the paint formula, the sale, the exhibitions, the juries, the reproductions, the dealer's third, the date the cheque cleared, and the corrections above the line, and the pencil afterthoughts made ten, twenty and thirty years later.

Even this needs a qualification. The transcription attributes hands cautiously, and it refuses to attribute the drawn lettering of the covers to either of them \q{until somebody has compared it against a signed drawing}. More importantly, Book IV, the cash book, is in a hand the transcriber could not place: the first batch of it is \q{a fluent, sloping copperplate, unlike either the small upright running hand or the drawn title lettering recorded from Books I to III}, and its entries begin in November 1913, eleven years before Josephine Nivison entered the account. Whoever kept the cash book in the illustration years, it was not the woman who later kept the work books, and the edition marks every passage of those early leaves as unidentified. The hand that begins the archive is a hand nobody has named.

Later hands are few and they are worth listing, because each is a moment when the record passed out of the household. \q{Given by Helen Hayes to Smith College} under The MacArthurs' Home (\lf{Book II}{77}{17681}); \q{John Clancy cited value for insurance \$15000 - 1964} under Office at Night (\lf{Book II}{79}{16494}); a heavier blunt pencil on leaf 46 of Book II recording a resale through Maynard Walker in 1956. These are the traces of the dealer, the buyer and the museum reading the book after the Hoppers had stopped writing in it.

\section{What a leaf is}

A leaf of the work books has a shape, and the shape is an argument about what a picture is. At the head, the title and the size in his hand. Below it the drawing. Under the drawing a line of paint: the colourman, the canvas and its priming, the white, the oil, the thinner. Then the description in her hand, which on the early leaves is an anecdote about the people in the picture and on the later leaves is a note of what colour everything is. Then the sale: buyer, price, \q{- 1/3 =} the net, the date. Then the exhibitions and their verdicts, the reproductions, the prizes, and whatever she thought of afterwards. A picture, on this evidence, is an object of a certain size, made of certain materials, showing certain people, that went to certain places and was paid for on a certain day. It is not an idea, and it is not a feeling, and the record is nearly silent about what its maker meant by it. That silence is the subject of half of what follows.

The etchings leaves of Book I are ruled differently, and the ruling matters. Their columns read \q{Date | accepted / Refused | Exhibitions | Sold to, and terms | Received}, so that the first column dates the exhibition and not the sale, and one line carries an impression's whole history from jury to cheque. A refusal is recorded as carefully as an acceptance, one letter in its own narrow column, and the National Academy of Design, the transcription's glossary observes, \q{is the body that refused him most often}. Book IV is ruled by the stationer into a date column, a description column and a money column divided into dollars and cents by two printed red rules, and nobody heads a column anywhere in its 159 leaves.

\section{How the reading was made}

The claim this note makes is stated here so that every chapter can be measured against it: what the edition adds to the study of these books is not a finding but a condition, that any sentence read from them can be checked against the photograph of the leaf it came from, on one screen, by anybody. The transcription this study is written from was made in September 2026 by a language model, in batches of twelve sheets, each batch a fresh reading with a header that says what the batch covers, what was hard about it and what was decided. It is a first pass, and it says so on every file. A critical apparatus marks what was read, what was guessed, what could not be read and whose hand it was; nothing illegible is guessed, because, as the edition's method page puts it, \q{a guessed figure is a sale that did not happen}. The transcribed text is the source of record; the tables, the counts and the timeline are derived from it by scripts whose rules are published, and every figure in this study can be traced to a row.

I say this not to praise the method but to fix the status of what follows. The quotations are accurate to a reading that may be wrong in detail and that a second reader is expected to correct. The figures are accurate to rules that are stated and that a second reader may reject. What cannot be wrong, because it does not rest on any reading, is the shape of the archive: two hands, ten books, fifty-four years, and a woman recording, in a voice that is unmistakably her own, the work of a man who would not talk about it.
